% =====================================================================
%  Harald Høffding, Claudius Wilkens: Erkendelsens Problem (review)
%  In: Det nittende Aarhundrede. Maanedsskrift for Literatur og Kritik, udgivet af Georg Brandes og Edvard Brandes.  [Aargang 1875--76] Oktbr.--Decbr. 1875 og Jan.--Marts 1876.  Oktober 1875, section „Literatur“, pp. 81-85.
%  Antiqua.
%
%  TRANSCRIPTION -- GENERATED; DO NOT EDIT BY HAND.  Høffding's four pieces of 1875-76 (Herbert Spencers Filosofi; the
%  review of Wilkens; Montaignes Betydning i Etikens Historie; Benedikt Spinoza) are built by ONE pipeline
%  (ebook-method/hoeffding-1875-76/) into a combined working file, which split.py cuts into four transcriptions.
%  Rebuild from the repo root:   sh ebook-method/hoeffding-1875-76/rebuild.sh
%
%  STATUS: COMPLETE (2026-10-15): the volume's title page and the whole text, pp. 81-85.
%
%  EDITION.  The only printing.  Text of record: the Internet Archive's own Scribe scan of the University of Illinois copy, detnittendeaarh1875unse_0 (500 ppi colour JP2); extract ~/bibliotek/Høffding, Harald/1875-76-nittende-aarhundrede-okt-mar-pp81-85.pdf (SCANS.tsv, sha256), made by ia_extract.py
%  (title leaf, Indhold leaves, the article).  SECOND COPY (voter only): Google's scan of the Harvard copy, nittendeaarhund01unkngoog (its ABBYY OCR).
%  The page range is read off the volume's Indhold („Literatur.  Wilkens: Erkendelsens Problem. (Anmeldt af Dr. phil. Høffding) ... 81“) and the scan: the review opens p. 81 under the section head „Literatur.“ and ends on p. 85, signed, above the review of Høyens Skrifter (not transcribed).
%
%  ACCURACY (2026-10-15; one collation for all four texts, 57 working pages, ~17,100 words).
%  (a) The DRAFT is our own tesseract reading (tessdata_best: dan for Det nittende Aarhundrede, script/Fraktur for Nær og
%      Fjern; 300 dpi) of the Archive's JP2 images, laid over the working PDF and turned into LaTeX by draft.py.  It was
%      chosen over the Archive's own OCR (tesseract 5.2/5.3, hOCR) by measurement against the Harvard copy's ABBYY OCR
%      (words unmatched: 3.0% against 9.5% in the Montaigne volume; the Archive reads aa as åa and long s as f).  The
%      Archive's OCR is the WITNESS: collated word by word (collate/check.py): 900 disagreements, 76 settled by check.py.
%  (b) Rules (rules.py: the Archive's å for a, long s read f, glyph faults against attested words, 326) and a THIRD reading
%      voted: the ABBYY OCR of ANOTHER COPY (Google's scan of the Harvard volume), by context and by word alignment
%      (rules.py harv-*, vote2.py hv-*), never on figures and never making a word the corpus does not attest.  AUDIT: 30
%      letter-changing rule/vote decisions read blind at the image: 0 contradicted (6 undetermined crops).
%  (c) The rest of Det nittende Aarhundrede read at the image BLIND (round b1: 210 crops on 12 sheets, five readers, 23/24
%      planted controls right; ~80 decided by hand from the readers' lines).  OPEN 0.
%  (d) Every FOOTNOTE read whole at the image (readers N1, N2: 26 notes on 23 pages), set from the reading; the note marks
%      checked at zooms (four the draft had misplaced, moved).  STRUCTURE read at the image (reader K): the title pages, the
%      heads, the endings and signatures, the verse on p. 116.
%  (e) Nær og Fjern (FRAKTUR, two columns, no second copy): each column PROOFREAD WHOLE at the image against the draft by
%      two readers (P1, P2; 3 of 4 planted errors found): 171 corrections (proofapply.py); then an unattested-word sweep
%      and a long-s / ø sweep, every doubtful word read at a zoom (proof-fixes.tsv, final-fixes.tsv).
%  (f) LETTERSPACING (\emph{} here) MEASURED at the image (spacemeasure.py, white runs between letters, threshold 5 px: 31
%      runs), every run judged at the image by reader K (13 kept in the body; the notes' spacing read by N1/N2).  Italics:
%      none found by stroke slant (check.py --italics).
%  (g) QUOTATION MARKS: the print quotes „...“ only (every quotation read at the image in b1, the notes and the structure);
%      both OCRs read „ as , or ,, and “ as ' ” " »: set by class (quotes.py), outside the notes.
%  (h) Sweeps on the collated text, every doubtful spot read at a zoom: unattested words (229 words in no other
%      transcription: 12 words mended, the misprints below confirmed); specks both OCRs had read as marks (specksweep.py,
%      marks the Harvard copy lacks mid-sentence: 16 removed).  The TRANSLATORS' odd spots (64) were read at zooms before
%      the translation was finished: 22 more faults mended (specks, ɔ: read as 2:, a running head and a footnote left in
%      the text, two paragraph breaks at page turns), 3 misprints confirmed (final-fixes.tsv).
%  (i) SAMPLED PAGE CHECK (2026-10-15, budget-limited to 3 pages): pp. 113 (Spencer), 217 (Montaigne) and one Nær og Fjern
%      column (p. 2, right), chosen at random (seed 20261015, one per stratum), transcribed BLIND at the image by one
%      reader (1,143 words) and compared word by word (pagecheck.py): 10 differences; decided at zooms: 3 reader slips;
%      7 errors of the collated text (2 words: „for“ for „før“, „skikkelig“ for the print's „stikkelig“; 5 marks: four
%      specks both OCRs had read as stops or quotes, one stop hidden by a planted control) = ~6 per 1000 words, all mended.
%      The specks led to specksweep.py over the whole of Det nittende Aarhundrede: 16 more found at zooms and removed.
%  PRINTED AS IS (misprints kept, marked at the line):
%    nothing recorded.
%  CONVENTIONS.  Page turns: % --- p. N --- and \opage{N} (the printed page); a word divided over a page turn is
%  split with %.  Footnotes as \footnote at their marks, numbered per page as printed (*), **)).  \emph{} = letterspacing.
%  The running heads and the sheet signatures are not transcribed.  The ø rule (CLAUDE.md).
% =====================================================================
\documentclass[12pt,a4paper,openany]{book}
\usepackage[utf8]{inputenc}
\usepackage[T1]{fontenc}
\usepackage{libertinus}
\usepackage{libertinust1math}
\usepackage[danish]{babel}
\usepackage[protrusion=true,expansion=true]{microtype}
\usepackage{setspace}
\usepackage[margin=1.3in]{geometry}
\usepackage{fancyhdr}
\usepackage{hyperref}
\hypersetup{pdftitle={Anmeldelse af Cl. Wilkens, Erkendelsens Problem}, pdfauthor={Harald Høffding}, hidelinks}
\pagestyle{fancy}
\fancyhf{}
\fancyhead[LE,RO]{\thepage}
\fancyhead[C]{\textit{H. Høffding: Anmeldelse af Cl. Wilkens, Erkendelsens Problem (1875)}}
\fancypagestyle{plain}{\fancyhf{}\fancyhead[LE,RO]{\thepage}}
\setstretch{1.3}
\usepackage{marginnote}
\usepackage{graphicx}
\DeclareUnicodeCharacter{0254}{\reflectbox{c}}   % ɔ (id est)
\makeatletter
\renewcommand{\thefootnote}{\ifcase\value{footnote}\or *\or **\or ***\else\arabic{footnote}\fi)}   % the print marks each page's notes *), **)
\makeatother
\newcommand{\opage}[1]{\marginnote{\footnotesize\textit{[#1]}}}
\newcommand{\signature}[1]{\par\smallskip\hfill #1\hspace*{2em}\par}
\pdfstringdefDisableCommands{\renewcommand{\footnote}[2][]{}\renewcommand{\opage}[1]{}}

\begin{document}

\begin{titlepage}
\begin{center}
{\LARGE\bfseries Det nittende Aarhundrede.}\\[1em]
{\small\bfseries \emph{Maanedsskrift}}\\[0.3em]
{\small for}\\[0.3em]
{\large\sffamily Literatur og Kritik.}\\[1.2em]
{\small Udgivet af}\\[0.3em]
{\large Georg Brandes {\small og} Edvard Brandes.}\\[0.6em]
\rule{3em}{0.4pt}\\[0.6em]
{\small\bfseries 1875. Oktbr.---Decbr. og 1876. Jan.---Marts.}\\
\vfill
% [printer's device: a ring with the monogram GB]
{\large\sffamily \emph{København}.}\\[0.4em]
{\small Forlagt af den Gyldendalske Boghandel (F. Hegel).}\\
{\small Trykt hos J. Jørgensen \& Co.}\\[0.4em]
{\small 1876.}
\end{center}
\end{titlepage}
\clearpage

% --- p. 81 ---
\setcounter{footnote}{0}%
\opage{81}\begin{center}{\large \emph{Literatur}.}\\[0.3em]\rule{3em}{0.4pt}\end{center}\medskip

{\small\noindent \emph{Claudius Wilkens}: Erkendelsens Problem. Et Forsøg paa positiv Kriticisme. København 1875.}\medskip

Forf. af det foreliggende Arbejde gør det meget vanskeligt
selv for sympatetiske Læsere at yde ham det Bifald,
hans Arbejde virkelig fortjener, fordi han i selve Skriftet
forud har optaget og tilkendt sig langt mere, end selv hine
Læsere kunne indrømme ham. Det har været hans Arbejde
til stor Skade. Anmelderen kender flere Eksempler
paa Læsere, som bleve afskrækkede fra at gøre sig bekendte
med Bogen paa Grund af Forordet. Filosofiske Undersøgelser
burde ogsaa, især i disse skeptiske Tider, og især naar man
som Forf. udtrykkelig proklamerer at staa paa et eksakt
Standpunkt, fremtræde med større Beskedenhed; det er baade
til Videnskabens og dens Dyrkeres Tarv.

Men dette maa ikke bortlede Opmærksomheden fra, at
Forf. arbejder i en Retning og paa et Omraade, som kan
siges at staa paa Dagsordenen. Lige siden den hegelske
Skoles Opløsning har den Overbevisning banet sig Vej, at
den idealistiske Spekulation egenlig med Urette betragtede
sig som fuldendende Konsekvens af Kants store filosofiske
Reform. Denne Spekulation var en naturlig Reaktion mod
Kants subjektive og revolutionære Standpunkt, men har langt
fra udtømt eller endog blot havt fuldstændig Øje for Betydningen
af hans Lære. Den nyere Naturvidenskab, særlig Sansefysiologien,
og den realistiske Skole i England have ogsaa
vakt Opmærksomheden for den store Betydning af Kants
Grundtanke i al dens Simpelhed. At vi ikke opfatte Tingene,
om de ere i sig selv, men som de fremstille sig for os paa
Grund af vor Organisation, derom hersker der nu stedse
mere Enighed, og det har ingen Forsker saa klart set som
% --- p. 82 ---
\setcounter{footnote}{0}%
\opage{82}Kant. Det var Problemet om Erkendelsens Oprindelse,
der førte Kant til denne Indsigt. Overfor den realistiske
Lære, at al Erkendelse, alt Bevidsthedsindhold udenfra prægedes
ind i Bevidstheden, stod den idealistiske Lære om de
medfødte Ideer, der, om de ikke fra først af laa færdige i
Sjælen, dog efterhaanden ligesom afdækkedes deres Svøb.
Kant afgjorde Striden saaledes, at vi ikke optage Noget i
vor Bevidsthed, uden at det undergaar en ved vor hele Organisation
betinget Omdannelse, men at der paa den anden
Side ikke er os Andet medfødt end selve denne Organisation.
Rum, Tid og Nødvendighed o. lign. Bestemmelser tilkomme
derfor kun Tingene i og for vor Erkendelse; men Rumsanskuelsen
ligger dog ikke færdig i os fra først af; ikke den,
men Grunden til den er os medfødt, ligesom Forestillingen
om den røde Farve ikke er os medfødt, hvorimod det har
sin Grund i vor medfødte Natur, at vi faa Fornemmelsen
af Rødt, naar vi se Zinnober.

Forf, kalder sin Teori Kriticisme, fordi han ligesom
Kant lærer, at vi ikke kende Tingene i sig selv. Men
han kalder den \emph{positiv Kriticisme}, fordi han mener at
kunne udfylde Kløften mellem Erkendelsen og Tingen i sig
selv. Thi kende vi end ikke Tingene i sig selv, saa kende
vi dog efter hans Mening deres Forhold, vi have „en absolut
Erkendelse af Tingenes Relationer“. Den kantiske Erkendelseslære
korrigeres ved, at der tillægges saa vel Erkendelsens
Former som dens Stof en subjektiv og en objektiv
Side“. --- Dette sidste Udtryk er uheldigt. Kant tillægger
ogsaa Erkendelsens Former „en objektiv Side“; at vi anskue
Tingene i Rummet, er i Følge ham ligesaa vel begrundet i
Tingene som i vor egen Natur; men vi vide ikke, hvad det
er i Tingene, der svarer til vor Rumsanskuelse. Naar vore
Anskuelses- og Tænkeformer udspringe af vor Natur, saa vi
aldrig kunne abstrahere fra dem, men se Alt under Forudsætning
af dem, havde Kant fuldstændig Ret til at slutte, at
vi ikke vide, hvad Tingene i sig selv ere. Naar Rummet er
en \emph{Anskuelsesform}, er det selvmodsigende at paastaa, at
Tingene i \emph{sig selv, ɔ: afset fra enhver Betragter} ere
i Rummet. Dersom Forf. havde kendt den i de senere Aar
herom i Tyskland førte Forhandling, vilde han næppe saa
tillidsfuldt have talt om Kants Fejlslutning her. Navnlig
har Hermann Cohen i et særdeles grundigt og skarpsindigt
Skrift („Kants Theorie der Erfahrung“, Berlin 1871) vist
Kants konsekvente Tankegang paa dette Punkt.

% --- p. 83 ---
\setcounter{footnote}{0}%
\opage{83}Den omtalte Selvmodsigelse træder klart frem i et af
Forf.s egne Eksempler. Have vi to Stokke, den ene dobbelt
saa stor som den anden, saa maa Forholdet mellem dem ---
siger Forf. --- jo dog være 1/2, hvad enten vi se paa dem
eller ikke, altsaa i sig selv. Men et Forhold eksisterer kun
under Forudsætning af en Sammenligning; naar der Ingen
er, som sammenligner, kan man altsaa heller ikke tale om
Forholdet 1/2. At det har sin Grund i Tingenes Væsen, at
vi opfatte dem i matematiske Forhold, nægter som sagt
Kant ikke. Men Forf. statuerer en \emph{absolut} Erkendelse af
Tingenes Relationer; Tingene selv opfatte vi efter hans Indrømmelse
kun ved Symboler, men Forholdet mellem Symbolerne
og Forholdet mellem Tingene skal være et og det
samme, saa at Forf. endog vil have det udtrykt ved en
matematisk Ligning. Kant kunde indrømme, at vore Anskuelse
svare til Tingenes Forhold, men ikke, at de ligne
disse, end sige, at de ere lig dem („gemäss“, ikke „ähnlich“,
end sige „gleich“).

Det er vanskeligt at danne sig en bestemt Forestilling
om, hvorledes Forf. egenlig vil have sin Lære om den
absolute Erkendelse af Tingenes Relationer forstaaet. Ved
Afhandlingens mundtlige Forsvar erklærede han det for sin
Mening, at Relationerne mellem Forestillingerne vare de
. samme som (=) Relationerne mellem Tingene. Men i Afhandlingen
selv udtaler han sig meget vaklende. Pag. 25 (kfr.
P. 118) hedder det, at vi vel kun faa et subjektivt Billede af
Verden, men et rigtigt Billede af Tilværelsens Relationer
som om Erkendelsen ikke kan være rigtig, fordi den er subjektiv!
Er det f. Eks. ikke rigtigt, at Zinnoberet er rødt,
og dog er dette en subjektiv Fornemmelse, thi Zinnoberet
er kun rødt for vore Øjne. --- Pag. 131, 177 og 182 derimod
kaldes Opfattelsen af Tingenes Relationer symbolsk,
hvilket er det samme som subjektiv; her har da Forf. glemt
sin Teori. --- Pag. 29 beraaber Forf. sig paa Positivismen,
som stemmende med hans Teori, og Pag. 117 hedder det:
„en forudsigende Erkendelse er absolut“. Her forveksler
Forf. en eksakt, ɔ: fuldkommen tydelig og kvantitativ bestemt
Erkendelse med en absolut Erkendelse i Betydning af
en saadan, der ikke er bunden til det Fænomenale, men angaar
det i sig Værende. Den eksakte Videnskab angaar kun
Fænomenerne, men giver os en fuldstændig tydelig Indsigt i
deres Relationer; Positivismen søger at samle alle eksakte
Videnskabers Resultater i et System; men heller ikke den
% --- p. 84 ---
\setcounter{footnote}{0}%
\opage{84}gaar ud over det Fænomenale. --- Denne Vaklen og Uklarhed
er en naturlig Følge af Teoriens selvmodsigende Karakter.

Forf.s „Korrektur af Kant“ er saaledes baade overflødig
og urigtig. Men han mener endnu at have en anden Port
aaben til Erkendelse af det Absolute, af Tingen i sig selv.
„Selviagttagelsen“, „Selvfølelsen“, „Livsfølelsen“ aabner os
en Bro til Tilværelsens Væsen. Vort eget Væsen fornemme
vi umiddelbart og have altsaa her ét Punkt af Tilværelsen
oplukket for os; vi vide, hvorledes vi selv ere til Mode ved
at være til, og slutte ad Analogiens Vej deraf, hvorledes
andre eksisterende Væsener ere det. Denne Appel til Følelsen
for at stoppe et Hul i Teorien varsler ikke godt.
Forfatteren har selv skitseret dette Punkt saa utydelig, at
det sikkert kun har været hans Mening at prøve sig frem.
Hvilken Tvetydighed ligger der ikke i en Sætning som denne:
„Den dybeste Forstaaelse er i Virkeligheden Følelsens“ (Pag.
140)! Forf. tager her, som en af hans Opponenter ved den
mundtlige Disputats gjorde opmærksom paa, „Forstaaelse“ i
en ganske anden Betydning end ellers i Afhandlingen; at
forstaa bliver nu ikke længer paa rationel Maade at begribe,
men paa digterisk Vis at sympatisere med: i denne Forstand
kan man nok sige med Forf., at vi forstaa Dyret bedre end
Stenen; men om Begreb er der saa ikke mere Tale. Det er
en Illusion at mene, at vi bedre skulde begribe det Aandelige
end det Legemlige, fordi hint er os umiddelbart givet i
vor Bevidsthed. Til en Begriben hører at se Sagen i dens
Forhold til alt Andet, men saa er netop hin Umiddelbarhed
en Hindring. De psykologiske Fænomener ere lige saa vel
Fænomener som de legemlige; og der eksisterer ingen særlig
Logik for dem.

Pag. 123 definerer Forf. det at „føle sin Væren“ ved
at „opfatte sin Værens Værd“ og erklærer, at en Væren, der
ikke besidder eller selv er Genstand for en saadan Følelse,
er „værdiløs, betydningsløs, ɔ: meningsløs“. Her vises vi
over til det praktiske Omraade. Moralsk og historisk interessere
vi os ganske vist kun for det, der har „Værdi og
Betydning“; men hvad borger os for, at dette er noget Absolut?
Den højeste Maalestok, vi kunne lægge til Grund ved
Værdbestemmelsen, er dog kun relativ, hentet fra den menneskelige
Naturs Krav og Opgaver, og det er en dristig Dogmatisme,
Forf. her styrter sig over i ved at identificere disse
med Tilværelsens Kærne, en Dogmatisme af lignende Art
som den, Kants „moralske Tro“ forledte saa mange Kantianere
% --- p. 85 ---
\setcounter{footnote}{0}%
\opage{85}til, og som den unge Schelling bekæmpede i sine interessante
„Briefe über Dogmatismus und Kriticismus“. Den nylig
citerede Sætning findes næsten med de samme Udtryk hos
Lotze (Mikrokosmus I, Pag. 407). Men det er værdt at
lægge Mærke til, med hvilken Forsigtighed denne klare og
skarpsindige Tænker gaar frem paa dette Punkt. Det er
ikke i Erkendelseslæren, han udtaler denne Sætning; det er
ikke engang i Metafysiken. Naar Metafysiken skal have videnskabelig
Karakter som rationel Udvikling af vor Verdensopfattelses
Grundbestemmelse, kan den ikke optage praktiske,
æstetisk-etiske Postulater. Derfor bør i Følge Lotze (Metafysik
Pag. 135) „Udtrykket det sande Værende, der kun er
af æstetisk Betydning [ɔ: udtrykker en Værdibestemmelse]
udelukkes fra den metafysiske Sprogbrug“. En saadan Varsomhed
iagttager Forf. ikke, uagtet han fægter (eller vil
fægte) under Kriticismens Banner.

Der er sikkert Punkter nok, hvor Kants Teori maa
korrigeres. Den nyere engelske Psykologi har aabnet en
Række af elementære Undersøgelser, som dels berigtige, dels
gennemføre, hvad Kant vilde. Den nyere Naturvidenskab og
den induktive Logik have fremdraget helt nye Problemer.
Men i det Hele og Store vil sikkert den gamle Königsberger
endnu i lange Tider staa med uanfægtet Storhed som den
nyere Tids centrale Tænker. Hans Ideer danne en fast
Kærne, om hvilken Alt, hvad den senere Udvikling paa de
forskellige videnskabelige Omraader har bragt for Dagen, vil
kunne slutte sig sammen.
\signature{\textbf{Harald Høffding.}}

\end{document}
